\section{The finite Newton--Schulz map}
\label{app:numerics}

The arms' derivation treats $\Nten$ as if it produced a polar factor. It does
not, and quantifying the gap turned out to matter for two separate reasons: it
bounds how far the geometric motivation can be trusted, and it exposed an
efficiency lever that is independent of the output side. Both measurements below
are CPU studies on authenticated tensors saved from the runs; neither is a
training result.

\subsection{How far ten steps is from a polar factor}

On 32 saved $512\times512$ matrices drawn from a mid-training update, we
compared $\Nten$'s output against a partial polar reference computed by SVD in
FP64, zeroing directions below a relative threshold. The map's output sits at a
mean cosine of $0.979$ to that reference with a mean relative error of $0.211$.

Two consequences. First, identities that hold for an exact polar factor ---
notably $PP^{\Tr}=I$ for a full-row-rank wide matrix, which would make a
post-orthogonalization left factor look redundant --- do not transfer, and the
arms that apply only the outer left factor are therefore not trivially
equivalent to the baseline. Second, the finite map is sensitive: all 32
matrices are numerically rank 512 under this convention, but the batch-wide
smallest singular values are around $2\times10^{-10}$, and that
near-deficiency is what makes the iteration amplify tiny input differences.

\subsection{A direction plateau at seven steps}

Varying the step budget $q\in\{3,5,7,10\}$ at fixed coefficients and
normalization gives the geometry and cost in Table~\ref{tab:nsbudget}.

\begin{table}[t]
\centering
\caption{Newton--Schulz step budget against direction geometry and nominal
cost, on saved mid-training tensors. ``Products'' counts matrix multiplications
per step times steps. Cosine and relative error are means against the partial
polar reference, under single-BF16-operand emulation matching the training
arithmetic. Seven steps matches ten on both geometry measures at 30\% fewer
products; three steps discards real direction information.}
\label{tab:nsbudget}
\footnotesize
\begin{tabular}{rrrr}
\toprule
Steps $q$ & Nominal products & Mean cosine to polar & Mean relative error \\
\midrule
3 & 9 & $0.774263$ & $0.652564$ \\
5 & 15 & $0.951345$ & $0.309694$ \\
7 & 21 & $\mathbf{0.980342}$ & $\mathbf{0.211339}$ \\
10 & 30 & $0.979295$ & $0.211032$ \\
\bottomrule
\end{tabular}
\end{table}

Seven steps is indistinguishable from ten on direction geometry --- marginally
better on cosine, marginally worse on relative error --- while using 30\% fewer
matrix products. Five steps is measurably worse but may still be a reasonable
trade. This is a genuine efficiency lever, and because Newton--Schulz is the
dominant per-matrix cost in \emph{every} arm including the baseline, it would
benefit all of them.

It must not be described as an implementation speedup. Five- and seven-step
directions differ from ten-step directions by relative errors of $0.411688$ and
$0.375523$, so changing the budget changes the optimizer, and any such arm needs
its own quality comparison. Section~\ref{sec:nonprogram} argues the right first
test is a three-point budget comparison on the baseline alone, not the
54-configuration sweep that crossing this with every geometry would require.

\subsection{What the checkpoint-continuation failure did and did not show}

One wave failed a gate requiring that training continuously to update 24 match
resuming from a checkpoint at update 22. The mismatch is real and large: at
update 24, 460 of 2{,}191 rank-specific tensors exceed
$\text{rtol}=10^{-5}$, $\text{atol}=10^{-6}$, across all eight ranks, while at
update 22 all 2{,}127 compared tensors are exactly equal. We spent
considerable effort localizing it, and the outcome is worth recording compactly
because it is a cautionary tale about where such effort goes.

The divergence localizes to the Newton--Schulz map's input sensitivity rather
than to any incorrect kernel. Gradients are bitwise exact at the divergence
point; root inputs are exact apart from process-local graph flags; a tiny
difference in the map's input is amplified from a relative Frobenius error of
$1.45\times10^{-6}$ to $0.0369$ in its output, reproducibly, on GPU and in CPU
emulation. We then tested three candidate fixes, and all three failed:
\begin{itemize}
\item \textbf{Higher operand precision.} Compensated products that retain
  low-order BF16 components reduce the output pair's relative error from
  $0.0368$ (single product) through $6.79\times10^{-5}$ (three products) to
  $3.68\times10^{-5}$ (six products), against FP32's $3.74\times10^{-5}$ and
  FP64's $3.66\times10^{-5}$. Every one of them still fails the original
  tolerance, including the FP64 reference, so operand precision alone is not the
  cause.
\item \textbf{Fewer steps.} Every step budget also fails, so the amplification
  is not merely an artifact of iterating ten times.
\item \textbf{A cancellation-free start.} Replacing a subtraction of nearby
  inverse-quarter powers in the map's initialization with an equivalent
  positive-pair rational form improves start accuracy by four orders of
  magnitude --- worst-group errors fall from $6.1\times10^{-5}$ to
  $8.0\times10^{-9}$, and all 32 matrices pass an FP64 gate that the original
  form fails --- and does \emph{not} fix the continuation, because both forms
  are already exact for the continuous-versus-resumed comparison.
\end{itemize}

The honest summary is that a finite, rounded, near-rank-deficient polynomial
iteration is sensitive enough that bitwise checkpoint continuation is the wrong
gate for it, and that chasing numerical reproducibility in this component
consumed effort that the experiments in Section~\ref{sec:program} would have
used better. We record it because the failure is instructive and because two
of its by-products --- the step-budget plateau above and the improved start
form --- are useful independently of the gate that motivated them. Neither has
been validated at full model scale.
